\documentclass[../../main.tex]{subfiles}

\graphicspath{{\subfix{images/}}}

\begin{document}

\chapter{Stadiul Actual al Sistemelor Inteligente de Supraveghere Video}
\label{ch:stadiul_actual}

\section{Evoluția sistemelor de supraveghere video}
\label{sec:evolutie}

Dezvoltarea sistemelor de securitate și supraveghere video a fost marcată de o evoluție tehnologică continuă și accelerată, transformând supravegherea dintr-un simplu instrument pasiv de înregistrare, într-un ecosistem proactiv, capabil de analiză și decizie în timp real. Această tranziție istorică s-a desfășurat în mai multe etape majore, fiecare etapă apărând ca răspuns la limitările tehnologice ale generației anterioare: trecerea de la sistemele analogice la cele bazate pe IP, integrarea algoritmilor de analiză a conținutului video (VCA) și, cel mai recent, revoluția adusă de rețelele neurale și deep learning.

Inițial, sistemele de supraveghere în circuit închis (CCTV) operau exclusiv analogic, bazându-se pe transmiterea semnalului video continuu prin cabluri coaxiale către unități locale de stocare, precum DVR-urile (Digital Video Recorder). Deși și-au dovedit utilitatea zeci de ani, aceste arhitecturi prezentau limitări structurale severe în ceea ce privește scalabilitatea infrastructurii, calitatea imaginilor captate și vulnerabilitatea echipamentelor fizice la sabotaj. Paradigma a fost complet schimbată odată cu adoptarea la scară largă a camerelor IP, care au permis integrarea echipamentelor video în infrastructurile rețelelor de date existente (LAN/WAN) \cite{ip_surveillance}. Această migrare a facilitat nu doar o îmbunătățire radicală a rezoluției (prin introducerea senzorilor HD și multi-megapixel), ci a oferit și capabilități de management centralizat și monitorizare la distanță. Cu toate acestea, literatura de specialitate subliniază că tranziția către mediul digital a generat noi provocări inginerești, în special legate de limitările lățimii de bandă (bandwidth) necesare pentru transmiterea fluxurilor video de înaltă rezoluție și de riscurile de securitate cibernetică inerente rețelelor conectate la internet \cite{ip_surveillance}.

În ciuda imaginilor de înaltă fidelitate oferite de camerele IP, eficiența opera-țională a sistemelor de supraveghere rămânea fundamental dependentă de resursa umană. Monitorizarea pasivă, în care operatorii umani trebuiau să urmărească simultan numeroase ecrane, era ineficientă și predispusă la erori majore din cauza oboselii și a pierderii concentrării. Pentru a depăși acest obstacol, industria a dezvoltat sisteme de tip Video Content Analysis (VCA), marcând tranziția de la o supraveghere pasivă la una activă, orientată spre detectarea evenimentelor \cite{vca_optimization}. VCA a dotat sistemele cu capacitatea de a alerta automat operatorii pe baza unor scenarii predefinite, precum încălcarea unui perimetru (line crossing), detectarea obiectelor abandonate, numărarea persoanelor sau identificarea comportamentelor neobișnuite (loitering) \cite{vca_response}. Aceste sisteme au reprezentat un progres major față de tehnologiile timpurii de detecție a mișcării (Video Motion Detection), care erau adesea declanșate fals de factori de mediu (vânt, schimbări de iluminare sau mișcarea vegetației). Totuși, algoritmii VCA clasici se bazau pe reguli stricte, programate manual, ceea ce îi făcea inflexibili și le limita performanța în medii vizuale complexe, foarte aglomerate sau dinamice.

Următorul salt tehnologic major, care a redefinit complet capabilitățile ecosistemului de supraveghere, a fost inserția inteligenței artificiale prin intermediul algoritmilor de Deep Learning. Spre deosebire de sistemele VCA tradiționale, care depindeau de „extragerea manuală a caracteristicilor” (manual feature extraction) realizată de programatori, modelele de tip Deep Learning---în special Rețelele Neurale Convoluționale (CNN)---au introdus capacitatea de a învăța autonom direct din seturi masive de date \cite{pandey_neural}. Această inovație a eliminat barierele algoritmilor bazați pe reguli, permițând sistemelor să clasifice obiectele cu o precizie excepțională, reușind să facă distincția clară, în timp real, între un om, un vehicul sau un animal.

Mai mult decât simpla clasificare și filtrare a alarmelor false, Deep Learning-ul a transformat modul în care datele video sunt utilizate. Sistemele moderne dotate cu rețele neurale nu mai acționează doar reactiv, ci dezvoltă capabilități predictive \cite{pandey_neural}. Prin analizarea a mii de ore de fluxuri video și identificarea unor tipare complexe de comportament, aceste sisteme pot anticipa potențiale incidente de securitate sau anomalii înainte de materializarea lor. Prin urmare, camera de supraveghere a evoluat dintr-un simplu senzor optic, într-un „nod inteligent” (edge device) capabil de procesare locală avansată, integrat într-un ecosistem informațional proactiv și automatizat.

\section{Soluții comerciale existente}
\label{sec:solutii_comerciale}

Piața sistemelor de supraveghere video este dominată în prezent de soluții comerciale robuste, dezvoltate de companii consacrate. Deși aceste platforme oferă capabilități avansate de inteligență artificială, ele funcționează ca sisteme închise (\emph{closed-source}), impunând adesea restricții legate de hardware și costuri recurente semnificative. În continuare sunt analizate trei dintre cele mai reprezentative produse din această categorie, pentru a evidenția limitările care motivează abordarea propusă în prezenta lucrare.

\subsubsection*{1. Verkada}
Verkada este o platformă modernă de securitate, nativă în cloud (SaaS - \emph{Software as a Service}), recunoscută pentru ușurința în instalare și managementul centralizat \cite{verkada_overview}.
\begin{itemize}
	\item \textbf{Capabilități AI declarate:} Sistemul integrează algoritmi de analiză video la nivel de procesor (\emph{edge computing}) capabili de detectarea persoanelor și a vehiculelor, recunoașterea atributelor faciale și a culorii hainelor, recunoașterea plăcuțelor de înmatriculare (LPR) și detectarea comportamentelor neobișnuite \cite{verkada_overview}.
	\item \textbf{Model de licențiere:} Utilizarea platformei necesită plata unei licențe software anuale pentru fiecare cameră în parte (cuprinsă adesea între 200 și 1800 USD per an/dispozitiv).
	\item \textbf{Limitări observate:} Cea mai mare constrângere este \emph{dependența totală de hardware-ul producătorului} (\emph{vendor lock-in}). Platforma software \linebreak Verkada funcționează exclusiv cu camerele produse de ei \cite{verkada_security}. În plus, costul total de proprietate (TCO - \emph{Total Cost of Ownership}) crește masiv pe termen lung din cauza licențelor recurente, iar dependența de cloud poate fi un dezavantaj pentru instituțiile care impun izolarea strictă a datelor.
\end{itemize}

\subsubsection*{2. Avigilon (Motorola Solutions)}
Avigilon este unul dintre liderii tradiționali ai pieței de sisteme de management video (VMS), oferind soluții hibride și \emph{on-premise} destinate mediului \emph{enterprise} \cite{avigilon_appearance}.
\begin{itemize}
	\item \textbf{Capabilități AI declarate:} Tehnologia sa de vârf, \emph{Avigilon Appearance Search}, permite căutarea rapidă a persoanelor și vehiculelor în zeci de fluxuri video simultan \cite{avigilon_docs}. Oferă, de asemenea, LPR, detecția mișcărilor neobișnuite (\emph{Unusual Motion Detection}) și integrare profundă cu sistemele de control al accesului.
	\item \textbf{Model de licențiere:} Licențierea este modulară, presupunând o achiziție inițială substanțială (per server și per canal/cameră), plus costuri suplimentare pentru modulele de analiză AI avansată și actualizări de mentenanță.
	\item \textbf{Limitări observate:} Deși software-ul poate accepta camere de la terți (prin protocolul ONVIF), obținerea performanțelor AI maxime și utilizarea anumitor algoritmi necesită achiziționarea fie a camerelor Avigilon cu AI integrat, fie a unor echipamente de procesare dedicate (\emph{Avigilon AI Appliances}) \cite{avigilon_appearance}. Sistemul funcționează ca o „cutie neagră” și presupune bugete inaccesibile utilizatorilor de rând sau organizațiilor mici.
\end{itemize}

\subsubsection*{3. BriefCam}
BriefCam este o soluție software pură, faimoasă pentru inovația sa în domeniul \emph{Video Synopsis} (rezumarea video), capabilă să suprapună evenimentele din momente diferite pe un singur fundal static pentru a condensa ore de înregistrări în câteva minute de vizionare \cite{briefcam_datasheet}.
\begin{itemize}
	\item \textbf{Capabilități AI declarate:} Oferă analiză video avansată în timp real și post-eveniment: extragerea datelor demografice (vârstă, gen), LPR, recunoaștere facială, clasificarea precisă a obiectelor și analiza comportamentală complexă (direcție, staționare, viteză) \cite{briefcam_datasheet}.
	\item \textbf{Model de licențiere:} Costurile sunt foarte ridicate, sistemul fiind destinat aplicațiilor critice (guvernamentale, aeroporturi, orașe inteligente). Licențierea se face pe pachete de bază la care se adaugă costuri per canal video procesat și per utilizator concurent.
	\item \textbf{Limitări observate:} Pe lângă modelul opac de procesare a datelor, algoritmii săi complecși impun o barieră hardware majoră. Sistemul necesită achiziționarea unor servere \emph{on-premise} extrem de puternice, echipate cu multiple plăci video dedicate (GPU) din gama enterprise, crescând masiv complexitatea instalării și a mentenanței \cite{briefcam_hardware}.
\end{itemize}

\subsubsection*{Concluzii privind soluțiile comerciale}
Deși produsele prezentate mai sus demonstrează nivelul ridicat de performanță la care a ajuns viziunea artificială aplicată în securitate, ele confirmă limitările identificate în Secțiunea~\ref{sec:motivatie}. Dependența de hardware proprietar (precum în cazul Verkada și, parțial, Avigilon), costurile uriașe de licențiere și lipsa de transparență a modului în care rețelele neuronale iau decizii justifică pe deplin nevoia unei alternative. Soluția dezvoltată în cadrul prezentei lucrări demonstrează că funcționalitățile esențiale — precum identificarea vehiculelor, detectarea armelor și analiza comportamentului — pot fi implementate utilizând exclusiv tehnologii cu sursă deschisă (\emph{open-source}), rezultând o arhitectură transparentă, adaptabilă și independentă de costuri de licențiere.



\section{Soluții open-source și academice}
\label{sec:solutii_open_source}

Alternativa la sistemele comerciale o reprezintă comunitatea cu sursă deschisă (\emph{open-source}) și cercetarea academică. Aceste inițiative elimină costurile de licențiere și dependența de hardware-ul proprietar, dar vin adesea cu limitări în ceea ce privește optimizarea pentru cazuri de utilizare foarte specifice sau lipsa unei arhitecturi software moderne, pregătite pentru utilizatori finali. Mai jos sunt analizate trei exemple relevante.

\subsubsection*{1. Frigate NVR}
Frigate este unul dintre cele mai populare sisteme de tip NVR (\emph{Network Video Recorder}) din comunitatea open-source, fiind construit nativ în jurul conceptului de inteligență artificială \cite{frigate_nvr}.
\begin{itemize}
	\item \textbf{Capabilități AI declarate:} Rulează inferență AI complet locală, evitând transmiterea datelor în cloud. Este capabil să analizeze fluxuri video în timp real și este optimizat pentru a folosi acceleratoare hardware precum Google Coral Edge TPU, Intel OpenVINO sau NVIDIA TensorRT \cite{frigate_nvr}. Oferă integrare nativă profundă cu platforme de automatizare precum Home Assistant.
	\item \textbf{Model de licențiere:} Complet gratuit și open-source (licență MIT).
	\item \textbf{Limitări față de soluția propusă:} Principala limitare a platformei constă în faptul că modelele AI furnizate implicit, antrenate în general pe seturi de date precum COCO, sunt concepute pentru detecția unor categorii de obiecte de interes general. Implementarea unor funcționalități specializate presupune integrarea și configurarea unor modele personalizate suplimentare. Astfel, în timp ce Frigate este orientat în principal către scenarii de monitorizare și automatizare specifice mediului \emph{smart home}, soluția propusă în această lucrare este dezvoltată pentru a răspunde cerințelor specifice aplicațiilor de securitate și supraveghere inteligentă.
\end{itemize}

\subsubsection*{2. ZoneMinder}
ZoneMinder este una dintre cele mai cunoscute platforme VMS (Video Management System) din ecosistemul open-source, fiind dezvoltată și întreținută în mod continuu de peste două decenii \cite{zoneminder_docs}.
\begin{itemize}
	\item \textbf{Capabilități AI declarate:} Platforma oferă suport pentru un număr mare de camere IP prin protocolul RTSP și include funcționalități avansate pentru gestionarea înregistrărilor și filtrarea evenimentelor pe baza unor zone definite de utilizator. În versiunile recente, au fost integrate și funcționalități de recunoaștere a obiectelor prin intermediul unor mecanisme externe bazate pe scripturi și \emph{webhook}-uri (de exemplu, servere de evenimente care interoghează scripturi Python) \cite{zoneminder_docs}.
	\item \textbf{Model de licențiere:} Gratuit și open-source (licență GPL).
	\item \textbf{Limitări față de soluția propusă:} ZoneMinder utilizează o arhitectură monolitică, implementată preponderent în C++ și PHP. Deși oferă un set extins de funcționalități, integrarea componentelor de inteligență artificială se realizează prin mecanisme suplimentare, externe fluxului principal de procesare. În multe configurații, analiza bazată pe AI este declanșată doar după detectarea unei modificări la nivelul imaginii (detecția clasică a mișcării), ceea ce poate introduce întârzieri suplimentare în procesarea evenimentelor. În contrast, sistemul propus în această lucrare, construit pe baza tehnologiilor FastAPI și React, adoptă o arhitectură separată pentru backend și frontend și efectuează inferența direct asupra fluxului video, reducând astfel latența procesării.
\end{itemize}

\subsubsection*{3. Implementări academice bazate pe YOLO (Custom Pipelines)}
O mare parte din literatura de specialitate modernă explorează utilizarea algoritmilor de detecție din familia YOLO (versiunile v5, v8, v10) specific pentru domeniul securității.
\begin{itemize}
	\item \textbf{Capabilități AI declarate:} Cercetările academice demonstrează că modelele YOLO antrenate pe seturi de date specifice ating o precizie excelentă (mAP, F1-Score) și un FPS ridicat în detecția armelor albe sau de foc, identificând amenințările în fracțiuni de secundă direct din imagini CCTV \cite{fierro2026,mandal2026}.
	\item \textbf{Model de licențiere:} Surse deschise (adesea scripturi publicate pe \linebreak GitHub în scop de cercetare) \cite{mandal2026}.
	\item \textbf{Limitări față de soluția propusă:} Majoritatea cercetărilor academice se opresc la stadiul de „Proof of Concept” (PoC). Soluțiile prezentate în lucrări științifice sunt, de regulă, simple scripturi izolate de Python care analizează videoclipuri preînregistrate sau cadre individuale \cite{fierro2026,mandal2026}. Acestora le lipsește complet un ecosistem aplicativ real: nu dețin un server backend care să gestioneze mai multe conexiuni RTSP, nu oferă o bază de date pentru stocarea și alertarea evenimentelor și nici o interfață grafică interactivă. Chiar și abordările academice mai avansate, care integrează un sistem complet de tip IoT cu notificări și control hardware \cite{khan2026}, rămân prototipuri proof-of-concept dedicate unui scenariu restrâns, fără a oferi un produs software generic end-to-end. Soluția propusă depășește acest stadiu teoretic, transformând modelele AI într-un produs software complet end-to-end.
\end{itemize}

Așa cum reiese din analiza realizată, soluțiile existente din ecosistemul open-source prezintă o serie de limitări în raport cu obiectivele urmărite în această lucrare. Unele platforme sunt concepute ca soluții generale de supraveghere video și nu sunt optimizate pentru scenarii specifice, în timp ce altele se bazează pe arhitecturi software dezvoltate pentru cerințe tehnologice anterioare sau sunt prezentate exclusiv în context experimental și de cercetare. Prin integrarea unor modele YOLO personalizate într-o stivă tehnologică modernă, bazată pe \linebreak FastAPI pentru serviciile backend și React pentru interfața utilizatorului, proiectul propus urmărește să combine performanța modelelor de detecție validate în literatura de specialitate cu cerințele de scalabilitate, modularitate și utilizare practică specifice unei aplicații software complete.

\section{Stadiul cercetării în domeniile relevante}
\label{sec:stadiul_cercetarii}

\section{Plasarea proiectului în context}
\label{sec:plasare_context}

\end{document}
